% =========================================================
\theme{T3}{Theoretical Exposition}
% =========================================================

\textbf{Core issue:} R2 (major a, d) and minor comments.
Clarify novelty vs.\ Camerlenghi et al.\ (2022),
add practical interpretation of results.

\textbf{Status: 9/10 items addressed} in current revision (red text).

\bigskip

\begin{longtable}{p{1cm} p{1.5cm} p{1cm} p{7cm} p{2.5cm}}
\toprule
\textbf{Item} & \textbf{Src} & \textbf{Pri} & \textbf{Description} & \textbf{Status} \\
\midrule
\endhead

T3.a & R2-a & --- & Clarify contribution vs.\ Camerlenghi et al. & \done \\
\midrule
T3.b & R2-d & --- & Practical interpretation after Cor.~3.5 and Prop.~3.7. & \done \\
\midrule
T3.c & R2-f & --- & Reorganize Section~3.1. & \done \\
\midrule
T3.d & R2-g & --- & Clarify NB parametrization. & \done \\
\midrule
T3.e & R2-h & --- & Clarify prior importance. & \done \\
\midrule
T3.f & R2-e & --- & First-trigger-time modeling advantage. & \done \\
\midrule
T3.g & R2-l & --- & Model selection guidance in Discussion. & \done \\
\midrule
T3.h & R2-b & --- & Clarify difference from algorithmic approaches. & \done \\
\midrule
T3.i & R2-i & --- & Fix SSP vs.\ Be-SSP naming. & \done \\
\midrule

T3.j & R2-l & \med &
  \textbf{Model fit assessment procedure.}
  R2 asks how to select the score distribution beyond guidance.
  Consider: marginal likelihood comparison, posterior predictive checks.
& \notstart \\

\bottomrule
\end{longtable}
